
\documentclass{beamer}

\usetheme{Madrid}

\title{Introduction to Technical Writing}
\author{Aakriti}
\institute{Computer Science and Engineering}
\date{\today}

\begin{document}

% Slide 1: Title
\begin{frame}
    \titlepage
\end{frame}

% Slide 2: Introduction
\begin{frame}{Introduction}
    Technical Writing is used to communicate
    technical information clearly and effectively.

    \begin{itemize}
        \item Reports
        \item Research papers
        \item Presentations
        \item Technical documentation
    \end{itemize}
\end{frame}

% Slide 3: LaTeX
\begin{frame}{LaTeX}
    LaTeX is a document preparation system.

    \textbf{Features:}
    \begin{itemize}
        \item Professional formatting
        \item Mathematical equations
        \item Tables and figures
        \item Automatic numbering
    \end{itemize}
\end{frame}

% Slide 4: Beamer
\begin{frame}{Beamer}
    Beamer is used to create presentations using LaTeX.

    \begin{itemize}
        \item Creates professional slides
        \item Supports different themes
        \item Includes equations and tables
    \end{itemize}
\end{frame}

% Slide 5: Doxygen
\begin{frame}{Doxygen}
    Doxygen is a documentation generation tool.

    \begin{itemize}
        \item Generates software documentation
        \item Uses special comments in source code
        \item Supports C++ and other languages
    \end{itemize}
\end{frame}

% Slide 6: Git and GitHub
\begin{frame}{Git and GitHub}
    \textbf{Git:}
    \begin{itemize}
        \item Tracks changes in files
        \item Manages project versions
    \end{itemize}

    \textbf{GitHub:}
    \begin{itemize}
        \item Hosts Git repositories
        \item Supports collaboration
    \end{itemize}
\end{frame}

% Slide 7: Conclusion
\begin{frame}{Conclusion}
    Technical Writing tools help us create
    professional documents and presentations.

    \begin{center}
        \textbf{Thank You!}
    \end{center}
\end{frame}

\end{document}
